\par\noindent\SourceLeaf{56}%
Diese mit dem vorangestellten Faktor multiplizierte Reihe lässt sich mittels der bekannten Formel

\[
\sin\varphi\cos\lambda\varphi=\frac{1}{2}\sin(\lambda+1)\varphi-\frac{1}{2}\sin(\lambda-1)\varphi,
\]

\noindent in Glieder zerlegen, die nach Sinus von Vielfachen fortschreiten. Dabei zerfällt jedes Glied außer dem Anfangsglied und dem ersten in zwei Glieder. Das eine erhält das nächstkleinere, das andere das nächstgrößere Vielfache als dasjenige, dessen Kosinus im zu zerlegenden Glied vorkommt. So liefert zum Beispiel das Glied

\[
\nu n . \sin\varphi\,A^{(2)}\cos2\varphi
\]

\noindent die beiden Glieder

\[
-\nu\frac{nA^{(2)}}{2}\sin\varphi+\nu\frac{nA^{(2)}}{2}\sin3\varphi;
\]

\noindent Fügt man das erste davon dem aus dem Anfangsglied hervorgehenden $\nu nA\sin\varphi$ hinzu, so erhält man das Anfangsglied der entwickelten Reihe

\[
=\left(\nu nA-\nu\frac{nA^{(2)}}{2}\right)\sin\varphi.
\]

Frage ich allgemein nach dem Koeffizienten, den $\sin(h+1)\varphi$ in der entwickelten Reihe erhält, so können von den Gliedern der gegebenen Reihe nur $\nu nA^{(h)}\sin\varphi\cos h\varphi$ (der erste Teil seiner Entwicklung ist

\[
+\nu\frac{nA^{(h)}}{2}\sin(h+1)\varphi)
\]

\noindent und $\nu n.A^{(h+2)}\sin\varphi\cos(h+2)\varphi$ (der zweite Teil seiner Entwicklung ist

\[
-\nu\frac{nA^{(h+2)}}{2}\sin(h+1)\varphi),
\]

\noindent Glieder liefern, die das gesuchte Vielfache von $\varphi$ enthalten. Somit ist in dieser Reihe der vollständige Koeffizient von $\sin(h+1)\varphi$

\[
=\nu\frac{n.A^{(h)}}{2}-\nu\frac{nA^{(h+2)}}{2}.
\]

\Needspace{8\baselineskip}
$\beta$) Auf der anderen Seite erhält man dagegen

\[
(1+n\cos\varphi)\left(A^{(1)}\sin\varphi+2A^{(2)}\sin2\varphi+.\,.\,.+hA^{(h)}\sin h\varphi+.\,.\,.\,.\,.\right)
\]

\noindent Bei der Entwicklung dieses Ausdrucks entstehen zwei Reihen. Die eine ist

\[
A^{(1)}\sin\varphi+2A^{(2)}\sin2\varphi+.\,.\,.\,.+(h+1)A^{(h+1)}\sin(h+1)\varphi+.\,.\,.\,.
\]

\noindent und ergibt sich unmittelbar.
